% Provenance: copied from the design workspace CambridgePQ/recursion.fast, `report/preamble.tex`, on 2026-10-19. Read-only copy; the source is authoritative.
% Shared preamble for the recursion.fast blueprint. Compile with: latexmk -lualatex recursion-fast-blueprint.tex
\usepackage[margin=2.4cm]{geometry}
\usepackage{fontspec}
\usepackage{amsmath,amssymb,amsthm,mathtools}
\usepackage{unicode-math}
\setmainfont{STIX Two Text}
\setmathfont{STIX Two Math}
\setmonofont{DejaVu Sans Mono}[Scale=0.82]
\usepackage{microtype}
\usepackage[section]{placeins}  % floats never drift past their section
\renewcommand{\topfraction}{0.85}\renewcommand{\bottomfraction}{0.6}\renewcommand{\textfraction}{0.12}
\renewcommand{\floatpagefraction}{0.75}\setcounter{topnumber}{3}\setcounter{totalnumber}{5}
\usepackage{booktabs,tabularx,longtable,multirow,array}
\usepackage{xcolor}
\usepackage{enumitem}
\usepackage[most]{tcolorbox}
\usepackage{fancyvrb}
\usepackage{tikz}
\usetikzlibrary{arrows.meta,positioning,fit,calc,shapes.geometric,shapes.multipart,backgrounds}
\usepackage[backend=biber,style=alphabetic,maxbibnames=99,maxalphanames=4]{biblatex}
\usepackage{csquotes}
\usepackage{hyperref}
\hypersetup{colorlinks=true,linkcolor=blue!50!black,citecolor=green!40!black,urlcolor=blue!60!black}
\usepackage[capitalise,noabbrev]{cleveref}

% ---------- theorem-like environments
\theoremstyle{definition}
\newtheorem{definition}{Definition}[section]
\newtheorem{example}[definition]{Example}
\theoremstyle{plain}
\newtheorem{theorem}[definition]{Theorem}
\newtheorem{lemma}[definition]{Lemma}
\newtheorem{proposition}[definition]{Proposition}
\theoremstyle{remark}
\newtheorem{remark}[definition]{Remark}


% ---------- boxes
% Decision box: \begin{decision}{D-ID}{Title} ... \end{decision}
\newtcolorbox{decision}[2]{enhanced,breakable,colback=blue!3,colframe=blue!45!black,
  fonttitle=\bfseries\small,title={Decision #1: #2},left=4pt,right=4pt,top=2pt,bottom=2pt}
% Issue box (actionable): \begin{issue}{I-ID}{Title} ... \end{issue}
\newtcolorbox{issue}[2]{enhanced,breakable,colback=red!3,colframe=red!50!black,
  fonttitle=\bfseries\small,title={Issue #1: #2},left=4pt,right=4pt,top=2pt,bottom=2pt}
% Trusted assumption box
\newtcolorbox{assumption}[2]{enhanced,breakable,colback=orange!4,colframe=orange!60!black,
  fonttitle=\bfseries\small,title={Assumption #1: #2},left=4pt,right=4pt,top=2pt,bottom=2pt}
% Key takeaway
\newtcolorbox{takeaway}[1][]{enhanced,breakable,colback=green!3,colframe=green!40!black,
  fonttitle=\bfseries\small,title={#1},left=4pt,right=4pt,top=2pt,bottom=2pt}

% Lean code blocks: use \begin{leancode} ... \end{leancode} (fancyvrb; Unicode-safe under LuaLaTeX).
% (listings is not used: it reorders non-ASCII characters under LuaLaTeX.)
\DefineVerbatimEnvironment{leancode}{Verbatim}{fontsize=\small,frame=single,framerule=0.3pt,
  rulecolor=\color{black!30},framesep=4pt,xleftmargin=2pt,xrightmargin=2pt}

% ---------- macros
\newcommand{\F}{\mathbb{F}}
\newcommand{\Fp}{\mathbb{F}_p}
\newcommand{\negl}{\mathsf{negl}}
\newcommand{\R}{\mathcal{R}}
\newcommand{\Rrcs}{\mathcal{R}_{\mathsf{R1CS}}}
\newcommand{\Racc}{\mathcal{R}_{\mathsf{acc}}}
\newcommand{\Prv}{\mathsf{P}}
\newcommand{\Vrf}{\mathsf{V}}
\newcommand{\Ext}{\mathsf{E}}
\newcommand{\Ind}{\mathsf{I}}
\newcommand{\vk}{\mathsf{vk}}
\newcommand{\idx}{\mathsf{IDX}}
\newcommand{\eps}{\varepsilon}
\newcommand{\errcap}{\varepsilon_{\max}}
\newcommand{\code}[1]{\texttt{#1}}
\newcommand{\lean}[1]{\texttt{#1}}
\newcommand{\repo}[1]{\texttt{#1}}
